\section{Development Toward Pose-Marginal Structural Validation}
\label{sec:mixturedevelopment}
This separate development changes the inference target from an arbitrary
continuous-density functional to compatibility of a specified structure with
held-out particles. It aims to remove supplied small pose balls. The current
calculations do not yet yield a useful continuous-pose test or experimental
calibration. Earlier failed approaches remain part of the evidence.

\paragraph{Established statistical construction.}
Let $Y=(Y_1,\ldots,Y_n)$ denote independent evaluation particles with known CTFs.
For a fixed candidate map $f$, let $p_{i,f,\nu}(Y_i\mid R)$ be a normalized
image likelihood at orientation $R$, where $\nu$ denotes other nuisance
parameters. A shared unknown viewing measure $G$ gives
\[
 p_{f,G,\nu}(Y)=\prod_{i=1}^n\int_{SO(3)}p_{i,f,\nu}(Y_i\mid R)\,dG(R).
\]
The viewing measure is assumed independent of the image-specific CTFs.
Let $q(Y)$ be a normalized predictive density learned independently of $Y$,
and let $D_f(Y)\ge\sup_{G,\nu}p_{f,G,\nu}(Y)$ pointwise. Then
$E_f(Y)=q(Y)/D_f(Y)$ has expectation at most one under every model in the
candidate null: integrate $q(Y)p_{f,G,\nu}(Y)/D_f(Y)\le q(Y)$.
Markov's inequality controls rejection at $E_f\ge1/\alpha$ by $\alpha$.
This is an application of established universal inference
\citep{universal2020}, not a new general coverage theorem. Normalization,
independence, and a genuine likelihood upper bound are indispensable.
An oracle predictive density checks algebra but does not demonstrate a learned
procedure. Noise misspecification or structural heterogeneity can invalidate
the null model. A rejected fixed structure does not identify the incorrect
feature or supply a confidence interval for arbitrary density.

\paragraph{Classical mixture dual.}
For strictly positive anchors $z_i$, any upper bound
$M\ge\sup_R\sum_i p_{i,f,\nu}(Y_i\mid R)/z_i$ yields
\begin{equation}\label{eq:mixturedual}
 \sup_G\log p_{f,G,\nu}(Y)
 \le \sum_i\log z_i+n\log(M/n).
\end{equation}
Apply arithmetic--geometric mean to the integrals divided by $z_i$ and
interchange finite summation and integration. This scaled likelihood dual,
finite-support maximizers and directional optimality checks are classical
\citep{lindsay1983}; scalable fixed-support optimization is also established
\citep{npmle2024}. A finite orientation grid alone supplies a feasible lower
likelihood, not a continuous upper. Nonuniform viewing laws must be treated
explicitly: recent misspecified-likelihood analysis and joint-estimation
consistency rely on stated identifiability and viewing-density conditions
\citep{nonuniformmle2025}.

\paragraph{Continuous Fourier Gaussian enclosures.}
Use the Hermitian Gaussian expansion already defined for reconstruction and
cover $SO(3)$ by Euler boxes for $R_z(a)R_y(b)R_z(c)$ on
$[0,2\pi]\times[0,\pi]\times[0,2\pi]$. For box half-widths $h_j$, write
$H=\sum_jh_j$. A frequency of norm $K$ moves at most
$2K\sin(\min(H,\pi)/2)$ from its center. For a Gaussian of width $s$ and
distance range $[l,u]$, gradient and Hessian norms are bounded by the maxima of
$r s^{-2}e^{-r^2/(2s^2)}$ and
$(1+r^2/s^2)s^{-2}e^{-r^2/(2s^2)}$ on that range. Both maxima occur at
$r=\operatorname{clip}(s,l,u)$. Summing absolute coefficients over conjugate
centers gives $B_1,B_2$. Taylor's formula bounds the per-image mean remainder by
\[
 b_i=\tfrac12H^2\big\| |C_i| (K^2B_2+KB_1)\big\|_2,
\]
where $C_i$ includes supplied noise whitening. If $r_i$ is the center residual
and $J_i$ its mean Jacobian, every vector $u$ gives the convex dual lower value
\begin{align*}
 d_i(u)&=u^\top r_i-\tfrac12\|u\|^2-\sum_jh_j|(J_i^\top u)_j|\\
 &\le\tfrac12\min_{|v_j|\le h_j}\|r_i-J_iv\|^2.
\end{align*}
Thus $\max(\sqrt{2\max(d_i(u),0)}-b_i,0)$ lower-bounds the nonlinear
residual norm. Its negative half square upper-bounds the Gaussian log kernel.
Coordinate-wise Gaussian intervals supply an additional enclosure; their
intersection retains high-frequency residual information lost by a single
image-wide ball. Full derivations and the separately tested second-order
log-kernel refinement are in the repository's versioned theory notes.

For each box $C$, let $U_{iC}$ upper-bound its likelihood kernel. Every viewing
law induces box masses, so the finite mixture with kernel matrix $U$ relaxes
the continuous likelihood. Applying Equation~\eqref{eq:mixturedual} to that
matrix yields a global upper bound. Center kernels independently give a
feasible lower bound. Bisection preserves coverage; each child can inherit
its parent's upper bound. A wall or split limit returns an unresolved bracket,
not a local optimizer labeled as a certified maximum. All proofs here concern
real arithmetic; ordinary floating-point checks are not interval arithmetic.

\paragraph{Discrete orientation and noise-scale screens.}
The first declared screen uses 64 allowed orientations, 128 CTFs and 440 real
Fourier coordinates per geometry, with sixteen independent noise repetitions.
The density generators are unit-$L^2$ constant-cell EMDB maps; alternatives
attenuate a locked 20\,\AA\ region by 50\% or 100\%, or remove the entire signal.
The predictive numerator is the true known-noise uniform mixture, so its
zero false-rejection count is an oracle algebra check. With supplied noise,
full local removal is detected in 14/16, 9/16 and 0/16 repetitions for
10028, 10049 and 10076. Half removal is never detected. Profiling an independent
noise scale per image/orientation loses all local detections.

Profiling one common positive scale preserves 14/16, 8/16 and 0/16 detections;
all zero-signal alternatives are rejected. This uses a global precision
bracket: for fixed residuals, the optimized log mixture of
$\exp(-r_{iC}^2t/2)$ is convex in precision $t$, so endpoint upper chords plus
the Gaussian $nd\log t/2$ term bound each interval. All 192 cases reach the
declared one-log-unit optimization gap, with maximum 0.937 and at most three
visited nodes. Total runtime is 72.83 seconds. These repeated outcomes use the
same observations as the supplied-scale screen; they are not independent
confirmatory evidence, continuous-orientation results or experimental power.

\paragraph{Continuous-orientation failures and local diagnosis.}
The next declared calculation generates 128 new Haar-oriented particles per
geometry from the independently fitted Fourier Gaussian pilot, normalized by
its analytic whole-space $L^2$ Gram norm. CTFs and Gaussian noise scales are
supplied, and shifts are zero. The optimizer receives no generating poses.
All three computations exhaust 8,192 splits, leaving 8,224 boxes. Their
feasible-center/global-upper gaps are 26,832, 20,548 and 19,874 log units, far
above the target of one. Solver times are 431, 380 and 391 seconds.

A declared follow-up fits the envelope-mixture weights more closely. Its
relaxation gaps are 0.0295, 0.00991 and 0.00977; nevertheless the continuous
uppers remain 20,253, 17,988 and 17,077 above the best recorded feasible values.
Those values also consider a uniform mixture on the generating poses,
evaluated only after fitting. This oracle support is a feasible viewing law,
not the true Haar integral or exact optimum. The envelope-mixture primal must
not be substituted for a continuous feasible lower value.

All 54 prescribed local curvature diagnostics complete. For boxes centered
at true simulated poses, with a sum of Euler half-widths of $0.1^\circ$,
median upper-to-locally-feasible gaps decrease from 0.0890/0.0897/0.0464 to
0.00166/0.00205/0.00075. At $5^\circ$ and $10^\circ$ the retained minimum shows
no improvement. Nine local starts per case supply feasible values, not exact
box maxima. These results justify investigating more selective global
refinement, but establish neither a tight continuous denominator nor a useful
structural test. A practical independent predictive numerator, unknown shifts,
common-scale continuous optimization and acquisition-model validation remain
required. Likelihood-based structural comparison itself is longstanding prior
art, including BioEM and CryoLike; an integrated score cannot automatically
replace a normalized predictive density \citep{bioem2013,cryolike2025}.


The separately declared envelope-guided refinement adds 30,528, 32,000 and
32,768 splits, retaining complete covers of 38,752, 40,224 and 40,992 cells.
The first two hit their 900-second wall limits; the third hits its split limit.
Final upper-to-feasible gaps remain 15,892, 14,490 and 12,898 log units.
Runner time is 2,718 seconds. The last finite envelope fits themselves have
gaps 69.1, 52.1 and 89.7, distinct from the larger continuous gaps. No case
reaches the target. Historical anchor result labels use ``row max'' for the
scaled dual with row maxima as anchors; the independent-image upper is the
sum of those maxima, which can be larger.

A focused mathematical review suggests preserving per-frequency Taylor errors.
For complex-frequency residual blocks $r_q$, let their error disks have radii
$\epsilon_q=H^2|C_q|(K_q^2 B_{2q}+K_qB_{1q})/2$. Every realified vector $u$
gives the squared-residual lower bound
\begin{align*}
 D(u)={}&2u^\top r-\|u\|^2-2\sum_q\epsilon_q\|u_q\|\\
 &-2\sum_jh_j|(J^\top u)_j|.
\end{align*}
This follows by minimizing the elementary squared-norm Fenchel bound over
the product of disks and the Euler box. Independent conic primal checks and
dual replay verify implementation; no new statistical inference principle is
claimed. Retaining old bounds avoids worsening the upper when finite dual
iterates are suboptimal. The declared diagnostic repeats all 54 local boxes
and evaluates 95 prescribed coarse cells per geometry, retaining all outcomes.
At two degrees the first two local median gaps decrease slightly, from 23.26
to 22.92 and 29.84 to 27.71. No coarse-cell improvement occurs on any of the
36,480 image/cell coordinates, and all three cover uppers are unchanged.
Thus neither refinement establishes a practical continuous likelihood test.
